\begin{abstract}%摘要+摘要+摘要+摘要+摘要


	%第一段——问题重述+简要思想：首先简要叙述所给问题的背景和动机，并分别分析每个小问题的特点（以下以三个问题为例）。根据这些特点说出自己的思想：针对于问题1，采用。。。。。。。。的方法解决；针对问题2用。。。。。。。。的方法解决；针对问题3用。。。。。。。。的方法解决。
	我们基于题面给定的DWTS官方赛季周度数据表，构建了一条端到端、可复现的建模流水线。该流水线先将原始宽表规范化为两张主线数据集（周度面板与赛季特征），再在统一口径下生成可追溯的CSV/PDF产物并落盘至\texttt{outputs/}目录，保证“论文叙述---模型输入---图表证据”一致。
	
	%\textbf{}：加粗	
	\textbf{问题1：}我们将观众投票视为不可观测潜变量，在每个\texttt{season$\times$week}上利用淘汰/退赛结果与评委信号进行约束，推断周内相对观众份额的后验分布。核心产物为后验汇总表\texttt{outputs/predictions/mcm2026c\_q1\_fan\_vote\_posterior\_summary.csv}与不确定性诊断表\texttt{outputs/tables/mcm2026c\_q1\_uncertainty\_summary.csv}，并配套生成可直接用于论文的矢量图。
	%第二段——模型建立及求解结果：介绍思想和模型： 对于问题1我们首先建立了。。。。。。。。模型I。首先利用。。。。。。，其次计算了。。。。。。，并借助。。。。。。数学算法和。。。。。。软件得出了。。。。。。结论。
	
	
	
	
	\textbf{问题2：}我们在Q1后验输出基础上，对Percent与Rank两类投票汇总机制进行反事实对比，量化机制差异在赛季与周尺度的分布，并定位“高分歧周/赛季”。赛季层汇总结果写入\texttt{outputs/tables/mcm2026c\_q2\_mechanism\_comparison.csv}，并辅以可视化诊断图支持逐周核查。
	%（第3段）	对于问题2我们用。。。。。。。。
	
	
	
	\textbf{问题3：}我们在同一数据口径下将比赛表现分解为评委线（技术表现）与观众线（人气强度），并拟合可解释的混合效应/OLS模型以分析年龄、行业、地区、舞伴等因素的影响方向与区间。关键系数与置信区间汇总于\texttt{outputs/tables/mcm2026c\_q3\_impact\_analysis\_coeffs.csv}。
	
	\textbf{问题4：}我们在受控压力测试下评估多种新赛制候选机制（多档离群冲击强度），并同时报告技术保护、公平性、人气表达与鲁棒性等多目标指标。主线指标产物为\texttt{outputs/tables/mcm2026c\_q4\_new\_system\_metrics.csv}，并通过决策视图图表将多目标权衡转化为可执行的机制建议。
	%（第4段）	对于问题3我们用。。。。。。。。（模型的建立与求解结果的陈述中，思想、模型、软件和结果必须描述清晰，亮点详细说明需突出。针对不同问题可独立成段也可采用一段式仅用分号“；”分割，摘要只接受文字描述形式，不接受图表等其他方式）
	
	
	
	
	
	%（第5段）	优化结果及总结：在。。。。。。条件下，针对。。。。。。模型进行适当修改与优化，这种条件的改变可能来自你的一种猜想或建议。要注意合理性。此推广模型可以不深入研究，也可以没有具体结果。
	
	%注：字数300~600之间，需控制在一页；摘要中必须将具体方法、模型和所得结果写出来；摘要要求“总分总”，段开头可用“针对问题1，针对问题2，针对问题3..”或者“首先，然后，其次，最后”等词语进行有逻辑的论述。摘要是重中之重，必须严格执行！
	
	
	% 美赛论文中无需注明关键字。若您一定要使用，
	% 请将以下两行的注释号 '%' 去除，以使其生效
	\vspace{5pt}  %mm	毫米	1 mm = 2.845 pt   pt 点	1 pt = 0.351 mm
	\noindent\textbf{Keywords}: keyword 1,  keyword 2,  keyword 3,  keyword 4,  keyword 5
	%\keywords:关键词；\quad：空格
	%使用到的模型名称、方法名称、特别是亮点一定要在关键字里出现，3~5个较合适,用分号隔开
\end{abstract}
\maketitle  % 生成 Summary Sheet
